%%%%%%%%%%%%%%%%%%%%%%%%%%%%%%%%%%%%%%%%%%%%%%%%%%%%%%%%%%%%%%
%%%%%%%%%%%%% MA-XXXX Matemática Exploratoria %%%%%%%%%%%%%
%%%%%%%%%%%%%%%%%%%%%%%%%%%%%%%%%%%%%%%%%%%%%%%%%%%%%%%%%%%%%%
\newpage
\subsection*{MA-0152 Matemática Exploratoria}
\addcontentsline{toc}{subsection}{MA-0152 Matemática Exploratoria}
\hypertarget{MA0152}{}
\begin{refsection}
\begin{table}[H]
\centering{
\begin{tabular}{|ll|}
\hline
%\multicolumn{2}{|c|}{\textbf{Nivel de virtualidad del curso:} Virtual}\\ \hline
\textbf{Año:} I  & \textbf{Requisitos:} Ninguno \\ 
\textbf{Ciclo:} I  & \textbf{Correquisitos:} Ninguno \\ 
\textbf{Tipo de curso:} Teórico & \textbf{Horas presenciales por semana: 5}   \\ 
\textbf{Créditos:} 3 & \textbf{Horas de trabajo independiente por semana: 4}  \\ \hline
\end{tabular}
}
\end{table}

\subsubsection*{Descripción del curso}

Este es un curso de primer ciclo para la carrera de Matemática,
pensado como una puerta de entrada al quehacer matemático. En lugar
de partir de teoría abstracta, el curso propone al estudiantado
problemas concretos --principalmente de teoría de números,
combinatoria y geometría-- y lo acompaña en el proceso real de hacer
matemática: explorar casos particulares, detectar patrones, formular
conjeturas, y finalmente justificar por qué esas conjeturas son
ciertas (o encontrar un contraejemplo que las refute). La
representación visual y geométrica de los problemas se usa como una
herramienta central de este proceso, no solo como un tema más: dibujar
y visualizar un problema antes de resolverlo es, junto con buscar
patrones o probar casos particulares, una de las formas principales
en que el estudiantado explora y conjetura a lo largo del curso.

Ese proceso de observar un patrón y conjeturar es el razonamiento
\textbf{inductivo}; el proceso de justificar la conjetura a partir de
verdades ya establecidas es el razonamiento \textbf{deductivo}. El
curso trabaja ambos de manera intuitiva y operacional: no se exige el
formalismo de un curso de demostraciones, pero sí se incorporan, en el
camino y sin descuidar el punto de vista intuitivo y de exploración,
conceptos, definiciones y notaciones propias de las matemáticas de
niveles superiores --conjuntos, funciones, grafos, propiedades de los
enteros-- y, hacia el final del curso, la lectura y escritura de
afirmaciones con los cuantificadores $\forall$ y $\exists$, que son la
base de las pruebas que el estudiantado deberá construir en cursos
posteriores de la carrera.

Más que memorizar técnicas, el objetivo es que el estudiantado
experimente qué se siente hacer matemática, y salga del curso con la
confianza y las herramientas mínimas --incluyendo la intuición
geométrica y visual-- para enfrentarse a un problema nuevo.

Además, se busca fomentar una apreciación más profunda de la belleza
y la elegancia de las matemáticas a través de la comprensión de sus
fundamentos lógicos y de sus representaciones visuales.


% Este es un curso de primer ciclo para la carrera de Matemática, pensado como una puerta de entrada al quehacer matemático. En lugar
% de partir de teoría abstracta, el curso propone al estudiantado problemas concretos (principalmente de teoría de números, geometría y combinatoria) y lo acompaña en el proceso real de hacer matemática:
% explorar casos particulares, detectar patrones, formular conjeturas, y finalmente justificar por qué esas conjeturas son ciertas (o
% encontrar un contraejemplo que las refute).

% Ese proceso de observar un patrón y conjeturar es el razonamiento
% \textbf{inductivo}; el proceso de justificar la conjetura a partir de
% verdades ya establecidas es el razonamiento \textbf{deductivo}. El
% curso trabaja ambos de manera intuitiva y operacional, no se exige el formalismo de un curso de demostraciones pero sí introduce, en el camino, el vocabulario y la notación que el estudiantado
% necesitará en cursos posteriores: conjuntos, funciones, grafos,
% propiedades de los enteros, y --hacia el final del curso-- la lectura
% y escritura de afirmaciones con los cuantificadores $\forall$ y
% $\exists$, que son la base de las pruebas que verán en Cálculo
% (límites por definición) y en cursos posteriores de la carrera.

% Más que memorizar técnicas, el objetivo es que el estudiantado
% experimente qué se siente hacer matemática, y salga del curso con la
% confianza y las herramientas mínimas para enfrentarse a un problema
% nuevo.

% Más específicamente, el razonamiento inductivo es el proceso de sacar una conclusión general al observar un patrón basado en instancias específicas. Esta conclusión se llama hipótesis o conjetura. En la matemática, el razonamiento inductivo muchas veces lleva al descubrimiento de nuevos resultados, los cuales luego se tratan de demostrar por medio del método deductivo, en el cual se usan razonamientos lógicos para establecer nuevas verdades a partir de verdades previamente establecidas. 

% El enfoque principal de este curso radica en permitir que el estudiantado descubra el mundo de las matemáticas a través de la resolución de problemas, lo que a su vez actúa como un vínculo hacia áreas de las matemáticas más avanzadas. En particular, se realizarán exploraciones con objetos básicos de naturaleza discreta como teoría de números y combinatoria. También se utilizar\'an los contenidos para dar ejemplos iniciales de ciertas nociones de pruebas sencillas de carácter operacional. 

% Aunque se comienza desde un nivel básico, a lo largo del recorrido se incorporan conceptos, definiciones y notaciones propias de las matemáticas de niveles superiores, como es el caso de conjuntos y funciones, grafos, propiedades de los números enteros, por mencionar algunos, todo esto sin descuidar un punto de vista intuitivo y de exploración.

% Además, se busca fomentar una apreciación más profunda de la belleza y la elegancia de las matemáticas a través de la comprensión de sus fundamentos lógicos.




%Este curso es una transición entre el nivel de matemática que ha adquirido el estudiantado durante sus años de enseñanza media y un nivel de matemática superior, requerido en cursos posteriores. Entendemos matemática como conceptos y entes abstractos, sus propiedades y las relaciones entre estos. Dada la relevancia de esta comprensión, el curso se ubica en primer ciclo de carrera, acompañado por el curso MA-0100 Introducción al Cálculo, conformando el primer acercamiento a un nivel de matemática más avanzado.

%El curso ha sido diseñado para iniciar este acercamiento, retomando aprendizajes previos con el fin de reforzar y estructurar el desarrollo de habilidades operacionales. Estos conocimientos brindarán herramientas para el óptimo desempeño en cursos posteriores. Este curso se complementará con MA-2XXX Matemática Introductoria II.

%Estimular (fomentar) en el estudiantado capacidades de análisis e interpretación que permitan el desarrollo de habilidades operacionales y el primer estímulo de pensamiento abstracto, ambos conceptos piedras angulares presentes en los mecanismos de razonamiento posteriores.


\subsubsection*{Objetivos}
\textbf{General:} Iniciar al estudiantado en el quehacer matemático mediante la resolución de problemas de teoría de números, combinatoria y geometría, a través de la exploración, la conjetura y
la justificación de resultados.


\textbf{Específicos.} Durante el curso se espera que el estudiantado sea capaz de:

\begin{enumerate}
\item Comprender y apreciar la importancia de las demostraciones en matemáticas como herramientas fundamentales para establecer la veracidad de afirmaciones matemáticas.
\item  Aplicar el razomiento inductivo para establecer conjeturas o conclusiones a partir de patrones observados en instancias específicas. 
\item  Justificar la veracidad o la falsedad de propiedades referentes a patrones y caracter\'isticas de objetos matem\'aticos aplicando habilidades operacionales.
\item Reproducir procedimientos con los contenidos matemáticos del curso que introduzcan al estudiantado en la interpretación, el descubrimiento, la inferencia y hasta la predicción de patrones y conexiones dentro del ámbito matemático.
 %item Estimular (fomentar) en el estudiantado capacidades de análisis e interpretación que permitan el desarrollo de habilidades operacionales y el primer estímulo de pensamiento abstracto, ambos conceptos piedras angulares presentes en los mecanismos de razonamiento posteriores.
 %\item Justificar la falsedad de un enunciado proporcionando contraejemplos.
 \item Analizar objetos matemáticos a   través de sus propiedades básicas y sus relaciones con otras propiedades. 
 \item Desarrollar intuición geométrica y visual como herramienta para explorar, conjeturar y comunicar resultados matemáticos.
\item Aplicar el enfoque de trabajo colaborativo y la participación en discusiones matemáticas como métodos para abordar problemas, desarrollar intuiciones y consolidar nociones y conceptos.
\item Comunicar ideas matemáticas de manera clara y efectiva en forma oral y escrita.
\end{enumerate}

\subsubsection*{Contenidos}

\begin{enumerate}
\item Introducción a los modos de trabajo y desarrollo de la matemática: 
    \begin{enumerate}
        \item Contexto histórico del desarrollo de las matemáticas y su formalización.
        \item Pruebas matemáticas, ¿por qué las necesitamos?, ¿cómo las podemos encontrar y ¿cómo se pueden escribir?
     \item Razonamiento inductivo y razonamiento deductivo.
    \end{enumerate}

\item Estrategias generales de resolución de problemas. Ver por ejemplo \cite{Pol14} y \cite[Capítulo 6]{Gri18}

\item Resolución de problemas concretos. Los siguientes contenidos deben incluirse en las sesiones enfocadas en la resolución de problemas y así lograr el cumplimiento de los objetivos del curso.
    \begin{enumerate}
        \item Teoría de números elemental.
        \item Teoría de conjuntos.
        \item Relaciones y funciones.
        \item Combinatoria enumerativa.
    \end{enumerate}

Los siguientes se sugieren como opciones recomendables para complementar el desarrollo de los objetivos del curso.
    \begin{enumerate}
        \item Grafos.
        \item Conteo y cardinalidad básica. El principio del palomar.
        \item Algebra elemental.
        \item Geometría: euclideana, análitica, vectorial.
    \end{enumerate}

\item  Introducción a la formalización de cuantificadores.
  \begin{enumerate}
    \item El cuantificador universal ($\forall$) y el existencial
      ($\exists$): sintaxis básica y significado intuitivo, usando
      ejemplos ya trabajados en el curso (p.\ ej.\ ``todo entero
      par\ldots'', ``existe un primo que\ldots'').
    \item Combinación y orden de cuantificadores: por qué
      $\forall x\, \exists y\, P(x,y)$ no es lo mismo que
      $\exists y\, \forall x\, P(x,y)$, con ejemplos concretos donde
      el orden cambia el significado.
    \item Negación de proposiciones cuantificadas (leyes de De Morgan
      para cuantificadores): cómo se niega $\forall x\, P(x)$ y
      $\exists x\, P(x)$, y por qué esto importa para pruebas por
      contradicción o por contraejemplo.
    \item Estrategias de prueba según el cuantificador: cómo se
      inicia una prueba de $\forall x\, P(x)$ (tomar $x$ arbitrario)
      y cómo se inicia una de $\exists x\, P(x)$ (exhibir un testigo
      o construirlo).
    \item Práctica de demostraciones con cuantificadores anidados
      (tipo $\forall x\, \exists y\, P(x,y)$), usando afirmaciones
      de teoría de números y combinatoria ya vistas en el curso
      --por ejemplo, ``para todo entero $n$, existe un entero $m$ tal
      que\ldots''-- de modo que el estudiantado practique el orden en
      que se manejan las variables (quién se fija primero, quién se
      construye después) y llegue a cursos posteriores, como Cálculo,
      ya familiarizado con la estructura de este tipo de
      demostración, aunque en este curso no se trabaje contenido de
      cálculo.
  \end{enumerate}

     
%    \item Los n\'umeros naturales y enteros: problemas de conteo, combinatoria. Conjuntos finitos, evitando formalismos pero una buena introducci\'on de problemas para solidificar nociones para el curso siguiente: usar combinatoria para contar elementos, n\'umeros de funciones, etc etc. Divisibilidad, teorema de divisi\'on eucl\'idea. 
%    \item Inducci\'on matem\'atica y recursividad: principio de inducci\'on como m\'etodo de demostraci\'on. Aplicaciones a problemas de divisibilidad. Introducci\'on de sumatorias y productos. La suma geom\'etrica. Sumas telesc\'opicas. Sucesiones recursivas. 
%    \item Geometr\'ia: problemas de geometr\'ia anal\'itica. Parametrizaciones de rectas, planos? c\'onicas, transformaciones. 
%    \item N\'umeros complejos: motivaci\'on hist\'orica, definiciones, norma, forma polar, f\'ormula de Moivre, potencias, ra\'ices.
%     \item Funciones: Estudiar desde el punto de vista explorativo la noción de función y criterio, sobreyectividad, inyectividad, imagen directa e imagen inversa de conjuntos, pero desde el punto de vista de los ejemplos. 
\end{enumerate}



\subsubsection*{Metodología}

En este curso los contenidos se abordan de manera transversal como herramientas para el desarrollo de habilidades matemáticas como el razonamiento inductivo y deductivo por medio de un enfoque basado en la resolución de problemas. Para esto se consideran las siguientes pautas y sugerencias.


\noindent \textbf{Lineamientos metodológicos:} La persona docente a cargo de este curso debe respetar las siguientes pautas:
\begin{enumerate}
    \item Este no es un curso formal de demostraciones, se espera que el enfoque sea intuitivo y operacional, basado en un enfoque de resolución de problemas. (Ver por ejemplo los libros \cite{Gri18}, \cite{MS17}, \cite{RT94} o el clásico \cite{Pol14})
    \item No es imprescindible que las justificaciones adopten la forma de demostraciones formales, sin embargo es importante que cuenten con argumentos intuitivos, precisos y acertados.
    \item Utilizar herramientas tecnológicas y software especializado para aspectos de cálculo, visualización, conjeturas y exploración. No se espera el desarrollo de habilidades de programación, pero sí una familiarización con algunas herramientas. 
    \item Debe complementarse con el uso de herramientas de inteligencia artificial acotando aspectos éticos, morales así como limitaciones técnicas de la herramienta.
\end{enumerate}



\textbf{Sugerencias metodológicas:} Adicionalmente, se tienen las siguientes recomendaciones para la persona docente a cargo de este curso:
\begin{enumerate}
    \item Realizar trabajos grupales en donde se fomente la participación activa por parte del estudiantado.
    \item Poner los contenidos en un contexto hist\'orico que reflejen su desarrollo natural. 
    \item Se recomienda que las evaluaciones sean enfocadas en proyectos, trabajos en grupo y exposiciones. 
    \item Usar la representación geométrica o visual como estrategia por
  defecto al abordar los temas obligatorios del punto 3 de Contenidos
  (teoría de números, conjuntos, relaciones y funciones,
  combinatoria), no solo cuando se llegue al tema opcional de
  geometría. En particular:
  \begin{itemize}
    \item Tratar de hacer un diagrama o dibujo del problema como una estrategia de resolución de problemas más, con el mismo peso que buscar un patrón o resolver un caso más simple.
    \item Incluir, para cada tema obligatorio, al menos un problema cuya solución se apoye en una representación visual natural (por ejemplo: diagramas de Venn para conjuntos, arreglos de
      puntos para identidades numéricas o de conteo, el triángulo de Pascal para coeficientes binomiales, gráficas para relaciones y
      funciones).
    \item Usar las herramientas tecnológicas mencionadas en la sugerencia anterior no solo para calcular, sino para que el estudiantado visualice un problema antes de intentar    resolverlo.
  \end{itemize}
    
    \item Para la sección de formalización de cuantificadores (ver
  Contenidos, punto 4), se recomienda:
  \begin{itemize}
    \item Dedicarle entre 2 y 3 sesiones, una vez que el estudiantado
      ya haya trabajado suficientes resultados de teoría de números
      y combinatoria como para tener ejemplos propios con qué
      practicar.
    \item Retomar conjeturas o resultados que el propio grupo haya
      formulado o demostrado antes en el curso, y pedirles que los
      reescriban usando $\forall$ y $\exists$ explícitamente, en vez
      de introducir ejemplos nuevos y ajenos al curso.
    \item Priorizar la mecánica de la prueba --quién se fija primero
      y quién se construye después-- por encima de la notación
      formal: no se espera que el estudiantado memorice reglas de
      lógica proposicional, sino que reconozca el patrón ``fijo un
      elemento arbitrario\ldots luego debo construir o exhibir
      otro''.
    \item No convertir esta sección en un curso de lógica proposicional: basta con
      $\forall$, $\exists$, su negación, disyunción, conjunción y el orden en cuantificadores
      anidados.
    \item Hacer ejercicios en donde el estudiantado escriba una afirmación sencilla primero en palabras propias y luego la
      formalice con cuantificadores, junto con su prueba completa.
  \end{itemize}
    
\end{enumerate}

\subsection*{Criterios de evaluación}

Dado que este curso no exige el formalismo de un curso de
demostraciones, la evaluación de proyectos, trabajos en grupo y
exposiciones (ver Sugerencias metodológicas) debe centrarse en el
\emph{proceso} de hacer matemática, y no únicamente en la corrección
formal del resultado final. En particular, se sugiere valorar:

\begin{itemize}
  \item \textbf{Exploración:} evidencia de haber probado casos
    particulares o ejemplos concretos antes de formular una
    conjetura, más que la sola presentación de la conjetura ya
    formada.
  \item \textbf{Conjetura:} claridad al formular la conjetura y
    coherencia entre esta y los casos explorados que la motivaron.
  \item \textbf{Justificación:} un argumento intuitivo, preciso y
    acertado de por qué la conjetura es cierta (o de por qué se
    descarta mediante un contraejemplo), sin exigir el formalismo de
    una demostración rigurosa, pero sí que el argumento sea
    convincente y esté libre de errores lógicos evidentes.
  \item \textbf{Comunicación:} claridad al explicar, oralmente o por
    escrito, tanto el proceso de exploración como la justificación
    final, de forma que otra persona del curso pueda seguir el
    razonamiento.
  \item \textbf{Trabajo colaborativo:} participación activa de cada
    integrante del grupo en la discusión y construcción de la
    solución, no solo en la presentación final.
\end{itemize}

No se recomienda evaluar el rigor formal de las justificaciones como
si se tratara de un curso de demostraciones: un argumento intuitivo
pero correcto y bien comunicado debe valorarse igual o más que uno
con notación formal pero sin comprensión real del porqué.


\nocite{Eng93}
\nocite{Gri18}
\nocite{MS17}
\nocite{Pol14}
\nocite{RT94}
\nocite{Sti22}
\nocite{Tal13}






\printbibliography[heading=subbibliography]
\end{refsection}
